% !TEX program = xelatex
\documentclass[11pt,a4paper]{article}
\usepackage[margin=2.2cm]{geometry}
\usepackage{fontspec,amsmath,mathtools,unicode-math}
\setmainfont{FiraSans}[Extension=.otf,UprightFont=*-Book,BoldFont=*-SemiBold,
  ItalicFont=*-BookItalic,BoldItalicFont=*-SemiBoldItalic]
\setmathfont{Fira Math}
\setmathrm{FiraSans}[Extension=.otf,UprightFont=*-Book,BoldFont=*-SemiBold]
\usepackage[ukrainian]{babel}
\usepackage{graphicx,booktabs,array,xcolor,titlesec,tikz}
\usetikzlibrary{arrows.meta,positioning}
\usepackage[most]{tcolorbox}
\usepackage[hidelinks]{hyperref}

\definecolor{Blue}{HTML}{1F5C93}
\definecolor{Frame}{HTML}{B8CFD9}
\definecolor{Back}{HTML}{F2F2F2}
\definecolor{Mark}{HTML}{CCE3F9}
\titleformat{\section}{\Large\bfseries\color{Blue}}{\thesection.}{0.6em}{}
\titleformat{\subsection}{\large\bfseries\color{Blue}}{\thesubsection.}{0.6em}{}
\tcbset{lec/.style={enhanced,breakable,colback=Back,colframe=Frame,boxrule=2pt,arc=3pt,
  left=8pt,right=8pt,top=5pt,bottom=5pt}}
\newcounter{defn}\newcounter{thm}
\newtcolorbox{defbox}{lec,before upper={\stepcounter{defn}\textbf{Означення \thedefn.}\ }}
\newtcolorbox{thmbox}[1][]{lec,before upper={\stepcounter{thm}\textbf{Теорема \thethm.}\ #1 }}
\newtcolorbox{factbox}{lec,before upper={\textbf{Факт.}\ }}
\newcommand{\term}[1]{\colorbox{Mark}{\textit{#1}}}
\graphicspath{{../figures/}}
% автоматично згенеровано experiments/run_all.py
\newcommand{\lamOne}{0.9058}
\newcommand{\lamTwo}{-0.0003}
\newcommand{\lamThree}{-14.572}
\newcommand{\lamOneErr}{0.0011}
\newcommand{\lamSum}{-13.6667}
\newcommand{\DKY}{2.0621}
\newcommand{\DTwo}{2.072}
\newcommand{\DTwoErr}{0.001}
\newcommand{\rHom}{13.9266}
\newcommand{\rAlo}{24.0578}
\newcommand{\rAhi}{24.0582}
\newcommand{\transGamma}{2.21}
\newcommand{\transGammaErr}{0.14}
\newcommand{\predSlope}{1.096}
\newcommand{\predLam}{0.913}
\newcommand{\predGain}{2.52}
\newcommand{\doublingTime}{0.765}
\newcommand{\lnsixLam}{1.681}
\newcommand{\lnsixNpos}{13}
\newcommand{\lnsixDKY}{27.0}
\newcommand{\lnsixDoubling}{0.41}
\newcommand{\lnsixKS}{10.14}
\newcommand{\enkfRmse}{0.50}
\newcommand{\enkfRmseX}{1.45}
\newcommand{\enkfFree}{10.4}
\newcommand{\orderRK}{4.17}
\newcommand{\orderEuler}{0.97}
\newcommand{\TnumA}{10.4}
\newcommand{\TnumB}{15.7}
\newcommand{\TnumC}{21.6}
\newcommand{\Tnumf}{11.2}
\newcommand{\TnumSlope}{5.04}
\newcommand{\TnumSlopeTh}{4.42}
\newcommand{\ftleNeg}{29}
\newcommand{\regimeMean}{2.32}
\newcommand{\regimeMax}{18}
\newcommand{\mapMinSlope}{1.11}
\newcommand{\butterflyT}{25.3}
\newcommand{\roundErr}{5.8\cdot 10^{-4}}

% автоматично згенеровано experiments/lecture_study.py
\newcommand{\redTable}{0.4 & 3.73\% & 0.3523 & 0.3636 \\
0.2 & 1.70\% & 0.1789 & 0.1818 \\
0.1 & 0.80\% & 0.0902 & 0.0909 \\
0.05 & 0.39\% & 0.0453 & 0.0455 \\}
\newcommand{\redErrSmall}{0.39}
\newcommand{\Lmax}{11.9}
\newcommand{\tBoundUseless}{5.13}
\newcommand{\tActual}{23.0}
\newcommand{\revSlope}{12.4}
\newcommand{\revLost}{2.25}
\newcommand{\avgXXA}{62.80}
\newcommand{\avgXXB}{62.81}
\newcommand{\avgXYA}{62.80}
\newcommand{\avgXYB}{62.81}
\newcommand{\avgBZA}{62.80}
\newcommand{\avgBZB}{62.80}
\newcommand{\avgYYA}{81.20}
\newcommand{\avgYYB}{81.19}
\newcommand{\avgRhsA}{81.21}
\newcommand{\avgRhsB}{81.20}
\newcommand{\avgZA}{23.55}
\newcommand{\avgZB}{23.55}

\newcommand{\fig}[3][1]{\begin{center}\includegraphics[width=#1\linewidth]{#2}\\[2pt]
  \small\itshape #3\end{center}}
\setlength{\parindent}{0pt}\setlength{\parskip}{6pt}

\begin{document}
\begin{center}
{\color{Blue}\bfseries\LARGE Динамічні системи}\\[6pt]
{\color{Blue}\Large Дивний атрактор Лоренца та детермінований хаос}\\[4pt]
Розгорнуте дослідження до семестрового проєкту\\[2pt]
Сучок Катерина, Шопська Анастасія, Бугір Єлизавета · 2026
\end{center}

\textbf{Головне питання.} Система задана трьома простими рівняннями без жодної випадковості.
Чому ж її майбутнє неможливо передбачити надовго? Виклад іде за етапами математичного моделювання.

\begin{center}
\begin{tikzpicture}[every node/.style={draw=Blue,fill=Mark!60,rounded corners,align=center,
  text width=5.4cm,minimum height=1.1cm},>={Stealth}]
\node (a) {Змістовна модель (§1)};
\node (b) [above=0.9cm of a] {Математична модель (§2)};
\node (c) [right=1.6cm of b] {Математичні та чисельні результати (§3–7)};
\node (d) [below=0.9cm of c] {Інтерпретація, верифікація і валідація (§8–9)};
\draw[->,thick,Blue] (a) -- (b); \draw[->,thick,Blue] (b) -- (c); \draw[->,thick,Blue] (c) -- (d);
\end{tikzpicture}
\end{center}

\section{Змістовна модель}
Шар рідини (або повітря) підігрівають знизу й охолоджують зверху. Поки перепад температур малий,
тепло передається теплопровідністю і рідина нерухома. Коли перепад перевищує поріг, виникають
конвективні вали: тепла рідина піднімається, холодна опускається. За ще більшого нагрівання вали
починають нерегулярно змінювати напрямок обертання.

Припущення: рух двовимірний; густина залежить від температури лише в члені з силою тяжіння
(наближення Бусінеска); із нескінченного набору просторових гармонік залишаємо три.

\section{Математична модель}
Функцію течії $\psi$ та відхилення температури $\theta$ від лінійного профілю наближають так
(Зальцман, 1962; Лоренц, 1963):
\[
\psi\propto X(t)\sin(\pi a x)\sin(\pi z),\qquad
\theta\propto Y(t)\cos(\pi a x)\sin(\pi z)-Z(t)\sin(2\pi z).
\]
Підставивши це в рівняння Бусінеска, відкинувши вищі гармоніки та змінивши масштаби, дістають
\begin{tcolorbox}[lec]
\[
\dot x=\sigma(y-x),\qquad \dot y=x(r-z)-y,\qquad \dot z=xy-\beta z .
\]
$x$ — інтенсивність конвекції; $y$ — різниця температур висхідного та низхідного потоків;
$z$ — відхилення вертикального профілю температури від лінійного.\\
$\sigma$ — число Прандтля; $r=\mathrm{Ra}/\mathrm{Ra}_c$ — відносне число Релея (нагрівання);
$\beta=4/(1+a^2)$ — геометричний параметр. Класичні значення: $\sigma=10$, $r=28$, $\beta=8/3$.
\end{tcolorbox}
Найменше критичне число Релея $\mathrm{Ra}_c=27\pi^4/4\approx657.5$ досягається при $a^2=1/2$,
звідки $\beta=8/3$.

\textbf{Тип моделі за класифікацією з лекції 1:} детерміністична, динамічна, неперервна,
нелінійна (добутки $xz$, $xy$), багатовимірна (вимірність 3), дисипативна; якісна для реальної
атмосфери й кількісна для конвекції лише поблизу порогу.

\begin{defbox}
Систему $\dot{\mathbf x}=\mathbf v(\mathbf x)$, $\mathbf x\in\mathbb R^n$, права частина якої не
залежить явно від часу, називаємо \term{автономною динамічною системою}. Точки, де
$\mathbf v(\mathbf x_*)=0$, — її \term{стани рівноваги}.
\end{defbox}

\section{Коректність моделі за Адамаром}
\subsection{Існування та єдиність}
\begin{thmbox}
Для будь-якої початкової умови задача Коші для системи Лоренца має єдиний розв'язок, визначений
для всіх $t\ge 0$.
\end{thmbox}
\textit{Доведення.} Права частина — многочлени, тобто неперервно диференційовна функція, тому
локально існує єдиний розв'язок (векторний аналог теореми з лекції 2). Лишається показати, що він
не втікає на нескінченність за скінченний час. Візьмемо $V=rx^2+\sigma y^2+\sigma(z-2r)^2$.
Уздовж розв'язків
\begin{align*}
\tfrac12\dot V&=rx\,\sigma(y-x)+\sigma y\,(rx-y-xz)+\sigma(z-2r)(xy-\beta z)\\
&=-\sigma\bigl(rx^2+y^2+\beta(z-r)^2-\beta r^2\bigr).
\end{align*}
Кубічні та мішані члени взаємно знищуються. Отже, $\dot V\le 2\sigma\beta r^2$, тобто $V$ росте
щонайбільше лінійно й залишається скінченною в кожен момент. \hfill$\mdlgwhtsquare$

\begin{factbox}
Поза еліпсоїдом $rx^2+y^2+\beta(z-r)^2=\beta r^2$ маємо $\dot V<0$. Тому всі траєкторії зрештою
потрапляють в обмежену \term{поглинаючу область} і не виходять з неї.
\end{factbox}

\textbf{Порівняння з лекцією 2.} У моделі біологічного вибуху $\dot x=kx^2$ права частина теж
квадратична, але система «живе» лише до $t_*=1/(kx_0)$. У системі Лоренца квадратичні члени не
додають «енергії» $V$, тому час її життя нескінченний. У прикладі з феніксами $\dot x=k\sqrt{x}$
єдиність порушується, бо $\sqrt x$ не диференційовна в нулі; тут такого бути не може.

\subsection{Стійкість щодо похибок у вхідних даних}
Нехай $\delta$ — різниця двох розв'язків. Оскільки поле квадратичне,
$\dot\delta=J(\bar{\mathbf x})\,\delta$, де $\bar{\mathbf x}$ — середина між розв'язками, а $J$ —
матриця Якобі. Звідси за нерівністю Гронуолла
\[
\|\delta(t)\|\le\|\delta(0)\|\exp\int_0^t\mu(s)\,ds,\qquad
\mu=\lambda_{\max}\bigl(\tfrac12(J+J^{T})\bigr).
\]
На \emph{скінченному} проміжку часу розв'язок неперервно залежить від початкових даних, тож задача
коректна за Адамаром. Але множник росте експоненційно: на обчисленій траєкторії $\mu$ сягає
\Lmax, і вже при $t\approx\tBoundUseless$ гарантована оцінка перевищує розмір атрактора.
Фактично відстань $10^{-8}$ виростає до 1 приблизно за $t\approx\tActual$.
\fig[0.92]{lec_gronwall}{Рис. 1. Дві траєкторії з початковою відстанню $10^{-8}$. Оцінку (помаранчева)
обчислено вздовж цих траєкторій; вона правильна, але дуже груба.}
\begin{tcolorbox}[lec]
\textbf{Висновок.} Хаос не суперечить коректності. Існування, єдиність і неперервна залежність є,
але \emph{рівномірної за часом} стійкості немає: потрібна точність початкових даних росте
експоненційно з горизонтом прогнозу.
\end{tcolorbox}

\section{Стани рівноваги та їхня стійкість}
З $\mathbf v=0$: $y=x$, $x(r-1-z)=0$, $x^2=\beta z$. Тому
\[
O=(0,0,0)\ \text{(завжди)},\qquad
C_\pm=\bigl(\pm\sqrt{\beta(r-1)},\ \pm\sqrt{\beta(r-1)},\ r-1\bigr)\ \text{(при $r>1$)}.
\]
\subsection{Глобальна стійкість нерухомої рідини при \texorpdfstring{$r<1$}{r<1}}
Для $V=x^2/\sigma+y^2+z^2$ виділенням повного квадрата маємо
\[
\tfrac12\dot V=-\Bigl(x-\tfrac{r+1}{2}y\Bigr)^{2}-\Bigl(1-\bigl(\tfrac{r+1}{2}\bigr)^{2}\Bigr)y^2-\beta z^2<0
\quad\text{для }(x,y,z)\ne0,\ 0\le r<1 .
\]
Отже, \emph{кожна} траєкторія прямує до $O$: конвекції немає.

\subsection{Стійкість \texorpdfstring{$C_\pm$}{C±}}
Характеристичне рівняння в точках $C_\pm$:
\[
\lambda^3+(\sigma+\beta+1)\lambda^2+\beta(\sigma+r)\lambda+2\beta\sigma(r-1)=0.
\]
За критерієм Рауса–Гурвіца всі корені мають від'ємні дійсні частини, якщо $a_2a_1>a_0$, тобто
\[
r<r_H=\frac{\sigma(\sigma+\beta+3)}{\sigma-\beta-1}=\frac{470}{19}\approx 24.74 .
\]

\section{Одновимірна система всередині системи Лоренца}
Біля $r=1$ два напрямки згасають швидко, а один — повільно. Позначимо $\mu=r-1$ і
$u=(x+\sigma y)/(1+\sigma)$. Тоді точно
\[
\dot u=\frac{\sigma}{\sigma+1}\,x\,(\mu-z).
\]
Після згасання швидких мод стан лежить на \term{центральному многовиді}, де в головному порядку
$x\approx u$, $z\approx u^2/\beta$. Дістаємо одновимірну динамічну систему з лекції 2:
\begin{tcolorbox}[lec]
\[
\dot u=v(u),\qquad v(u)=\frac{\sigma}{\sigma+1}\,u\Bigl(\mu-\frac{u^2}{\beta}\Bigr).
\]
\end{tcolorbox}

\subsection{Якісний аналіз за графіком \texorpdfstring{$v(u)$}{v(u)}}
Стани рівноваги: $u=0$ і, при $\mu>0$, $u_*=\pm\sqrt{\beta\mu}$ — це в точності $C_\pm$.
Розв'язки зростають, де $v>0$, і спадають, де $v<0$. Оскільки $\ddot u=v'(u)\,v(u)$, точки
перегину лежать на рівнях $v'(u)=0$, тобто $u=\pm\sqrt{\beta\mu/3}$.
\fig[0.95]{lec_phase_lines}{Рис. 2. Ліворуч — функція швидкостей, праворуч — інтегральні криві.
Зелені стани рівноваги стійкі, помаранчевий нестійкий; жовтий пунктир — рівні перегину.}
При переході $\mu$ через нуль стійкий стан $O$ втрачає стійкість і народжує два нових стійких:
це \term{біфуркація «вила»}. Фізично виникає конвекція з одним із двох напрямків обертання.

\subsection{Точний розв'язок: знову логістичне рівняння}
Заміна $w=u^2$ дає $\dot w=2u\dot u$, тобто
\[
\dot w=\frac{2\sigma}{\sigma+1}\,w\Bigl(\mu-\frac{w}{\beta}\Bigr),\qquad
w(t)=\frac{\mu\beta}{1+\bigl(\mu\beta/w_0-1\bigr)e^{-2\sigma\mu t/(\sigma+1)}} .
\]
Отже, квадрат інтенсивності конвекції задовольняє \emph{логістичне рівняння}. При $\mu=0$ маємо
\[u(t)=\frac{u_0}{\sqrt{1+2\sigma u_0^2t/(\beta(\sigma+1))}}:\]
згасання вже не експоненційне, а степеневе
(\term{критичне сповільнення}).

\subsection{Перевірка редукції}
\fig[0.9]{lec_reduction}{Рис. 3. Повна тривимірна система та одновимірна редукція.}
\begin{center}
\begin{tabular}{c>{\centering\arraybackslash}p{3.6cm}cc}\toprule
$\mu=r-1$ & найбільша відносна похибка & $\lambda_+$ в $O$ (3D) & $\sigma\mu/(\sigma+1)$ (1D)\\\midrule
\redTable
\bottomrule\end{tabular}
\end{center}
Похибка спадає приблизно пропорційно $\mu$: при $r=1.05$ вона дорівнює \redErrSmall\%.

\subsection{Чому для хаосу замало однієї змінної}
\begin{factbox}
Кожен розв'язок одновимірної системи $\dot u=v(u)$ з гладкою $v$ монотонний: за єдиністю він не
може перетнути стан рівноваги, тому знак $v$ уздовж нього сталий.
\end{factbox}
Тож у системах лекції 2 неможливі навіть коливання. На площині обмежена траєкторія, за теоремою
Пуанкаре–Бендіксона, прямує до рівноваги, циклу або з'єднань рівноваг. Хаос в автономній
неперервній системі можливий лише за вимірності щонайменше 3.

\section{Вирішальний параметр: нагрівання \texorpdfstring{$r$}{r}}
Так само як квота вилову $c$ змінює портрет популяції, параметр $r$ змінює долю системи Лоренца.
\begin{center}
\begin{tabular}{>{\raggedright}p{3.4cm}p{11.2cm}}\toprule
$0\le r<1$ & усі траєкторії прямують до $O$ (доведено вище)\\
$r=1$ & біфуркація «вила», народжуються $C_\pm$\\
$1<r<\rHom$ & траєкторія йде до одного з $C_\pm$\\
$r=\rHom$ & гомоклінічний вибух: народжуються нестійкі цикли\\
$\rHom<r<\rAlo$ & перехідний хаос, який зрештою закінчується в $C_\pm$\\
$\rAlo<r<24.74$ & співіснують дивний атрактор і стійкі $C_\pm$\\
$r>24.74$ & $C_\pm$ нестійкі; при $r=28$ спостерігається дивний атрактор\\\bottomrule
\end{tabular}
\end{center}
Значення $1$ і $24.74$ отримано аналітично, $\rHom$ і $\rAlo$ — чисельно (бісекція). При
значно більших $r$ є також вікна періодичності.

\section{Дивний атрактор і хаос}
\subsection{Стискання об'єму}
\begin{thmbox}[(Ліувілля, застосування.)]
Об'єм $\Omega(t)$ довільної області фазового простору змінюється за законом
$\dot\Omega=\int\operatorname{div}\mathbf v\,dV$.
\end{thmbox}
Тут $\operatorname{div}\mathbf v=-(\sigma+1+\beta)=-41/3$, тому
$\Omega(t)=\Omega(0)e^{-41t/3}$: це знову рівняння радіоактивного розпаду з лекції 1, із
«періодом піврозпаду» об'єму $\ln2/(41/3)\approx0.05$. Гранична множина має нульовий об'єм.

\begin{defbox}
\term{Атрактор} — компактна інваріантна множина, яка притягує всі траєкторії з деякого свого околу
і не розкладається на менші такі множини. Атрактор називаємо \term{дивним}, якщо він має
фрактальну структуру; рух на ньому хаотичний, якщо близькі траєкторії розбігаються експоненційно.
\end{defbox}
\fig[0.75]{pop_attractor}{Рис. 4. Атрактор Лоренца при класичних параметрах.}

\subsection{Показники Ляпунова}
\begin{defbox}
\term{Найбільшим показником Ляпунова} називаємо
\[\lambda_1=\lim_{t\to\infty}\lim_{\delta_0\to0}\frac1t\ln\frac{\|\delta(t)\|}{\|\delta_0\|}.\]
\end{defbox}
Алгоритмом Бенеттіна ($T=10^4$) отримано
$(\lambda_1,\lambda_2,\lambda_3)=(\lamOne,\ \lamTwo,\ \lamThree)$. Перевірка: сума дорівнює
\lamSum, тобто дивергенції. Розмірність Каплана–Йорке $D_{KY}=\DKY$; кореляційна розмірність
$D_2=\DTwo$. Атрактор «товщий» за поверхню, але не заповнює об'єм.

Те, що атрактор класичної системи Лоренца справді дивний, строго довів В. Такер (2002) за
допомогою обчислень із гарантованими оцінками похибок.

\section{Інтерпретація результатів}
\subsection{Горизонт передбачуваності}
Якщо прогноз корисний, доки $\|\delta\|<\Delta$, то
\[
T_p\approx\frac1{\lambda_1}\ln\frac{\Delta}{\delta_0}.
\]
Зменшення початкової похибки в 10 разів додає лише $\ln10/\lambda_1\approx\predGain$ одиниць часу.
Порівняйте з часом подвоєння з лекції 1: похибка подвоюється за
$\ln2/\lambda_1\approx\doublingTime$.
\fig[0.8]{lec_horizon}{Рис. 5. Горизонт прогнозу ($\Delta=1$). Нахил середнього \predSlope{} близький до $1/\lambda_1$.}

\subsection{Чи можна відновити минуле?}
Математично так: за єдиністю розв'язок продовжується назад. Практично — лише на короткий час,
бо стискання вперед означає розтягування назад. Експеримент: інтегруємо вперед на $T$, потім назад
на $T$ у подвійній точності. Похибка росте приблизно як $e^{\revSlope T}$ (орієнтир
$|\lambda_3|=14.57$), і вже при $T=\revLost$ початковий стан повністю втрачено.
\fig[0.85]{lec_reversibility}{Рис. 6. Похибка повернення в початковий стан.}
Це перегукується з прикладом теплопровідності з лекції 1: дисипація стирає інформацію про минуле.

\subsection{Детерміноване, але зі стабільною статистикою}
Для обмеженого розв'язку середнє за великий час від повної похідної прямує до нуля. З похідних
$x^2$, $z$ та $y^2+(z-r)^2$ отримуємо
\[
\langle x^2\rangle=\langle xy\rangle=\beta\langle z\rangle,\qquad
\langle y^2\rangle=\beta\bigl(r\langle z\rangle-\langle z^2\rangle\bigr).
\]
\begin{center}
\begin{tabular}{lcc}\toprule
Часове середнє ($T=2\cdot10^4$) & старт $(1,1,1)$ & старт $(-7,3,30)$\\\midrule
$\langle x^2\rangle$ & \avgXXA & \avgXXB\\
$\langle xy\rangle$ & \avgXYA & \avgXYB\\
$\beta\langle z\rangle$ & \avgBZA & \avgBZB\\
$\langle y^2\rangle$ & \avgYYA & \avgYYB\\
$\beta(r\langle z\rangle-\langle z^2\rangle)$ & \avgRhsA & \avgRhsB\\
$\langle z\rangle$ & \avgZA & \avgZB\\\bottomrule
\end{tabular}
\end{center}
Окрему траєкторію надовго передбачити не можна, а середні («клімат») практично не залежать від
початкової умови. Це відмінність прогнозу погоди від прогнозу клімату в межах цієї моделі.

\subsection{Що робити на практиці}
Замість одного прогнозу запускають \term{ансамбль} із трохи різними початковими умовами, а нові
вимірювання регулярно враховують (\term{асиміляція даних}). У нашому експерименті з фільтром
Калмана (EnKF) середньоквадратична похибка стану становить \enkfRmse{} при шумі вимірювань 2.

\section{Верифікація і валідація}
\textbf{Верифікація} (чи правильно розв'язано рівняння):
\begin{itemize}
\item спостережуваний порядок методів: Рунге–Кутти \orderRK{} (теорія 4), Ейлера \orderEuler{} (теорія 1);
\item сума показників Ляпунова збігається з дивергенцією;
\item тотожності для часових середніх виконуються (таблиця вище);
\item окрема траєкторія при великих $t$ залежить від кроку й арифметики, тому порівнюємо
статистичні характеристики, а не точки.
\end{itemize}
\textbf{Валідація} (чи ті рівняння розв'язано): три гармоніки добре описують конвекцію лише
поблизу порогу $r=1$. При $r=28$ справжня рідина збуджує багато інших мод, тому модель не є
прогнозом атмосфери. Її цінність якісна: вона показує, що нерегулярність і обмежена
передбачуваність виникають уже в простій детерміністичній системі. Ті самі рівняння точніше
описують, зокрема, водяне колесо Малкуса.

\section{Висновки}
\begin{enumerate}
\item Задача Коші для системи Лоренца має єдиний глобальний розв'язок, обмежений при $t\ge0$.
\item Вона неперервно залежить від даних на скінченному часі, але не рівномірно за часом.
\item Біля $r=1$ система зводиться до одновимірної, а $u^2$ задовольняє логістичне рівняння.
\item Хаос потребує щонайменше трьох змінних і нелінійності.
\item При $r=28$ окремий прогноз обмежений горизонтом $T_p$, тоді як статистика стійка.
\end{enumerate}

\textbf{Питання для самоперевірки.} (1) Знайдіть стани рівноваги редукції при $\mu=0.3$ та рівні
перегину. (2) Чому з $\operatorname{div}\mathbf v<0$ ще не випливає хаос? (3) На скільки зросте
горизонт прогнозу, якщо похибку зменшити в 1000 разів?

\begin{thebibliography}{9}\small
\bibitem{l} E. N. Lorenz. Deterministic nonperiodic flow. \textit{J. Atmos. Sci.} 20 (1963), 130–141.
\bibitem{s} B. Saltzman. Finite amplitude free convection as an initial value problem. \textit{J. Atmos. Sci.} 19 (1962), 329–341.
\bibitem{sp} C. Sparrow. \textit{The Lorenz Equations: Bifurcations, Chaos, and Strange Attractors}. Springer, 1982.
\bibitem{st} S. H. Strogatz. \textit{Nonlinear Dynamics and Chaos}. 2nd ed. Westview, 2015.
\bibitem{t} W. Tucker. A rigorous ODE solver and Smale's 14th problem. \textit{Found. Comput. Math.} 2 (2002), 53–117.
\bibitem{b} G. Benettin et al. Lyapunov characteristic exponents for smooth dynamical systems. \textit{Meccanica} 15 (1980), 9–30.
\end{thebibliography}
\end{document}
